\section*{Chapitre \ref{ch:fm:treasury-and-funding} --- Trésorerie et financement}

\begin{solution}{exo:fm:treasury-and-funding:1}
A~: $0.04\times120+0.25\times120+0.15\times60=43.8$~; B~: $0.10\times120=12.0$~; C~: $0.05\times120+0.15\times120+0.30\times80=48.0$ (en millions de dollars).
\end{solution}

\begin{solution}{exo:fm:treasury-and-funding:2}
5~millions de dollars (50~\%) sur la position longue~; 150~\% de la valeur de la position courte, soit 15~millions de dollars, y compris les 10~millions de dollars de produit de la vente.
\end{solution}

\begin{solution}{exo:fm:treasury-and-funding:3}
$38.1\times5~\%=1.91$ millions de dollars par an.
\end{solution}

\begin{solution}{exo:fm:treasury-and-funding:4}
Poser $x_{ib}=1$ pour un seul courtier et 0 ailleurs satisfait $\sum_bx_{ib}=1$ et les bornes, les variables auxiliaires prenant
leurs plus petites valeurs admissibles~; le programme minimise sur un ensemble qui contient ce point, si bien que son minimum ne peut pas dépasser le coût de ce point.
\end{solution}

\begin{solution}{exo:fm:treasury-and-funding:5}
Sa marge, 72.0~millions de dollars, est presque égale à celle de A, 71.55~millions de dollars, mais ses spreads de financement sont les plus faibles~: 2.39~millions de dollars par an, contre 3.30
millions de dollars pour A.
\end{solution}

\begin{solution}{exo:fm:treasury-and-funding:6}
Dix jours repoussent le jour où la trésorerie s'épuise, du jour 5 au jour 11, juste après les hausses de marge des courtiers pendant la première semaine du scénario, mais ces hausses
arrivent alors~; trente jours couvrent tout le stress et la société ne manque jamais de liquidités.
\end{solution}

\begin{solution}{exo:fm:treasury-and-funding:7}
Avec 50 ou 60~millions de dollars, le portefeuille optimisé tient tout le mois~; celui confié au courtier B épuise sa trésorerie le jour 4 avec 50~millions de dollars et le jour 5 avec 60
millions de dollars.
\end{solution}

\begin{solution}{exo:fm:treasury-and-funding:8}
Une marge plus faible aujourd'hui n'apporte pas davantage de liquidité en situation de stress~: un seul courtier peut relever sa \omterm{def:fm:treasury-and-funding:house}{marge maison} à volonté, et, sans second courtier ni \omterm{def:fm:treasury-and-funding:lockup}{immobilisation de marge},
l'exigence sur l'ensemble du portefeuille bouge d'un coup. Dans le chapitre, le portefeuille confié à un seul courtier épuise sa trésorerie deux jours plus tôt que le portefeuille diversifié et nécessite
29.8~millions de dollars de \omterm{def:fm:treasury-and-funding:cash}{coussin de liquidité} supplémentaires.
\end{solution}

\begin{solution}{pb:fm:treasury-and-funding:1}
\begin{enumerate}
\item Voir la \cref{def:fm:treasury-and-funding:cash}.
\item La marge initiale exigée par les chambres de compensation a augmenté d'environ 300~milliards de dollars au cours du mois, et les appels de marge de variation ont été bien supérieurs à ceux de février.
\item Voir les \cref{def:fm:treasury-and-funding:house,def:fm:treasury-and-funding:lockup}.
\item Regulation T~: marge initiale de 50~\% sur les titres de capital achetés sur marge, 150~\% pour les ventes à découvert~; FINRA 4210~: marge de maintenance de 25~\% sur les positions longues, ou marge de
portefeuille.
\item Voir la \cref{def:fm:treasury-and-funding:optim}.
\item Voir la \cref{prop:fm:treasury-and-funding:lp}.
\item Voir le \cref{tab:fm:treasury-and-funding:alloc}~: de 33.9~millions de dollars et 4.98~millions de dollars par an dans le cas optimisé, à 144.0~millions de dollars et 10.92
millions de dollars tout auprès de C.
\item L'optimisation équilibre la part de chaque courtier pour éviter la charge nette de A et les majorations pour concentration de C, et confie à B les positions que son taux brut forfaitaire tarifie à bon marché.
\item Les relations, la révision des conditions, le coût du transfert et la valeur d'un courtier qui soutient la société en situation de stress.
\item Voir la \cref{def:fm:treasury-and-funding:horizon}.
\item Cinq jours~: positions longues à $-2\%$, positions courtes à $-1\%$, indice à $-2\%$~; les courtiers doublent la \omterm{def:fm:treasury-and-funding:house}{marge maison} dès le jour 5. La semaine fait perdre 15.5~millions de dollars.
\item 43.0 et 71.8 millions de dollars sans \omterm{def:fm:treasury-and-funding:lockup}{immobilisation de marge}~; environ 36 et 37~millions de dollars avec~; la marge du courtier compensateur de 2.02 et 2.13~millions de dollars.
\item Optimisé~: jour 5, \omterm{def:fm:treasury-and-funding:cash}{coussin de liquidité} supplémentaire de 8.1~millions de dollars~; un seul courtier~: jour 3, 29.8~millions de dollars de plus~; avec une immobilisation de trente jours, aucun déficit.
\item Elle repousse le déficit à la fin du délai de préavis (jour 6 pour cinq jours, jour 11 pour dix).
\item Voir l'\cref{dat:fm:treasury-and-funding:fsb}.
\item Par la \cref{met:fm:treasury-and-funding:buffer}~: le déficit cumulé le plus important sur les scénarios de stress, sur le délai de préavis le plus long.
\item L'\omterm{def:fm:treasury-and-funding:lockup}{immobilisation de marge}~: l'optimisation fait économiser 1~million de dollars par an mais ne raccourcit rien~; sans immobilisation, le portefeuille optimisé épuise encore sa trésorerie le jour 5.
\item Des immobilisations d'au moins la durée du stress pour les conditions de \omterm{def:fm:treasury-and-funding:house}{marge maison} qui évoluent le plus~; libération de l'excédent sur demande~; préavis avant toute modification
du collatéral éligible.
\item 38.1~millions de dollars (53~\%) de marge économisée par rapport au courtier le moins cher pris seul~; sans immobilisation, la trésorerie tient jusqu'au jour 5 (jour 3 avec un seul courtier),
avec une immobilisation de trente jours, la trésorerie tient tout le mois.
\item Répartir le portefeuille entre plusieurs courtiers au moindre coût qui reste acceptable quand les conditions changent, avec des immobilisations couvrant l'horizon de stress~;
conserver un \omterm{def:fm:treasury-and-funding:cash}{coussin de liquidité} égal au plus grand déficit de la prévision de stress, communiqué chaque mois avec l'\omterm{def:fm:treasury-and-funding:horizon}{horizon de survie}.
\end{enumerate}
\end{solution}

\begin{solution}{iq:fm:treasury-and-funding:1}
La marge qu'exige un courtier selon ses propres règles~; elle tarifie le risque que le courtier fait peser sur le client (concentration, liquidité, crédit du client) au-delà
de ce que suppose le minimum réglementaire.
\iqlookfor{le risque propre du courtier et sa discrétion.}
\end{solution}

\begin{solution}{iq:fm:treasury-and-funding:2}
Le nombre de jours avant que la \omterm{def:fm:treasury-and-funding:cash}{trésorerie non grevée} ne passe pour la première fois sous zéro, après tous les appels et sorties en situation de stress~; à partir d'une prévision quotidienne du P\&L et
de l'exigence de chaque compte.
\iqlookfor{une prévision quotidienne, pas un ratio.}
\end{solution}

\begin{solution}{iq:fm:treasury-and-funding:3}
Minimiser le coût de la marge plus le financement sur les fractions $x_{ib}$ avec $\sum_bx_{ib}=1$~; linéariser la valeur absolue de l'exposition nette et les excédents de concentration
au moyen de variables auxiliaires~; un programme linéaire.
\iqlookfor{la convexité et la linéarisation.}
\end{solution}

\begin{solution}{iq:fm:treasury-and-funding:4}
Savoir si la convention l'autorise sans préavis, la nouvelle exigence face à la \omterm{def:fm:treasury-and-funding:cash}{trésorerie non grevée}, ce qui peut passer à d'autres courtiers, et quelles positions
réduire~; puis payer, réduire ou transférer, et communiquer l'\omterm{def:fm:treasury-and-funding:horizon}{horizon de survie}.
\iqlookfor{conditions, trésorerie et options, dans cet ordre.}
\end{solution}

\begin{solution}{iq:fm:treasury-and-funding:5}
Pour garder plusieurs courtiers engagés, acheter une \omterm{def:fm:treasury-and-funding:lockup}{immobilisation de marge} ou un financement confirmé, et conserver les positions là où elles peuvent être transférées ou financées en situation de stress.
\iqlookfor{le prix de la résilience.}
\end{solution}

\begin{solution}{iq:fm:treasury-and-funding:6}
Cartographier les comptes et leurs conditions~; construire des scénarios de P\&L, de hausses de \omterm{def:fm:treasury-and-funding:house}{marge maison}, d'excédents retenus et de mouvements de marge en chambre de compensation~; prévoir la trésorerie au jour le jour~;
communiquer l'\omterm{def:fm:treasury-and-funding:horizon}{horizon de survie} et le \omterm{def:fm:treasury-and-funding:cash}{coussin de liquidité} nécessaire, ainsi que les conditions qui les déterminent.
\iqlookfor{scénarios conjoints et prévision quotidienne.}
\end{solution}
